\documentclass[10pt,twocolumn]{article}
\usepackage{amsmath,amssymb,amsthm,mathtools}
\usepackage[T1]{fontenc}
\usepackage{lmodern}
\usepackage[hidelinks]{hyperref}
\hypersetup{
  pdftitle={Causally Ordered Security-State Theory (COSS)},
  pdfauthor={Nicola Gallo}
}
\usepackage{microtype}
\usepackage{booktabs}
\usepackage{array}
\usepackage[letterpaper,margin=0.72in]{geometry}
\setlength{\columnsep}{0.24in}
\emergencystretch=3em
\raggedbottom

\usepackage{stmaryrd}
\usepackage{enumitem}
\usepackage{url}
\usepackage{xspace}

% Remove end-of-proof square
\renewcommand{\qedsymbol}{}

% -----------------------------------------------------------------------------
% Model naming is intentionally centralized. Rename here if the final acronym
% changes after the literature/trademark review.
% -----------------------------------------------------------------------------
\newcommand{\Model}{\textnormal{\textsc{COSS}}\xspace}
\newcommand{\ModelFull}{Causally Ordered Security-State Theory\xspace}
\newcommand{\PIC}{\textsc{PIC}\xspace}
\newcommand{\PoR}{\mathrm{PoR}}
\newcommand{\PoC}{\mathrm{PoC}}
\newcommand{\sem}[1]{\llbracket #1 \rrbracket}
\newcommand{\powerset}{\mathcal{P}}
\newcommand{\state}{\Sigma}
\newcommand{\coord}{\sigma}
\newcommand{\st}{\operatorname{st}}
\newcommand{\FP}{\mathrm{FP}}
\newcommand{\EP}{\mathrm{EP}}
\newcommand{\CDAF}{\mathrm{CDAF\text{-}E}}
\newcommand{\Influence}{\mathrel{\rightsquigarrow_{\!\FP}}}
\newcommand{\CausedBy}{\operatorname{CausedBy}}
\newcommand{\Realize}{\operatorname{realize}}
\newcommand{\Exec}{\operatorname{exec}}
\newcommand{\IndAuth}{\operatorname{IndependentAuthority}}
\newcommand{\Accepted}{\operatorname{Accepted}}
\newcommand{\hb}{\prec_{c}}

\newtheorem{definition}{Definition}
\newtheorem{theorem}{Theorem}
\newtheorem{proposition}{Proposition}
\newtheorem{lemma}{Lemma}
\newtheorem{corollary}{Corollary}
\newtheorem{remark}{Remark}

\title{Full Provenance on the COSS Causal Axis:\\
An Event-Level Closure of Causal Attribution and\\
Confused-Deputy Exclusion}
\author{Nicola Gallo}
\date{27 September 2026}

\begin{document}
\maketitle

\begin{abstract}
Provenance Identity Continuity (PIC) was developed to answer a question that
possession-based authorization leaves open: not only which authority an actor
holds, but which execution occurrence may exercise it. PIC binds authority to
execution occurrences in a causal history, relates adjacent occurrences by
Proof of Relationship (PoR), and requires the carried authority to be
non-expansive along that history. This paper asks whether that structure can
serve as a general theory of security-state continuity, and uses it as the
reference frame from which existing security models are read.

We develop \ModelFull (\Model) as the generalization of PIC. An occurrence
carries a product state $\state_i=(\coord_i^{(j)})_{j\in J}$; each coordinate
keeps its own refinement order, while PoR remains the single causal axis.
Product Continuity then follows componentwise: ordinary causal transport
cannot make any carried security state more permissive. PIC is recovered as
the authority-only instance and relational PIC as the two-coordinate instance;
lattice information flow, typed capability transfer, and decentralized labels
can be placed as further coordinates without altering their semantics.

We then read Rajani, Garg, and Rezk's 2016 formal study of confused-deputy
attacks through this frame. Their full provenance semantics tracks explicit
and implicit influence, and their Theorem~6 proves CDA-freedom with endorsement
unconditionally for their region calculus. Within \Model, that influence is one
instance of an influence relation on the same causal order used by PoR: we give
a conservative event refinement whose erasure leaves their semantics
unchanged, and we attribute to their work neither PoR, PoC, nor distributed
lineage. The results of this paper rely only on three closure properties of the
influence relation, of which full provenance is one realization.

Once all security-relevant effects and influence edges of the modeled machine
are represented as events, causal attribution is derived by reachability
inside the model. Combined with PIC authority continuity and effect admission,
this yields a request-bounded effect theorem: in a closed \Model execution,
an effect caused by a request can exercise only authority present in that
request's origin, or introduced by an explicit authorized composition or
relaxation event on its causal path, and the lineage-form confused-deputy
condition is unsatisfiable along ordinary continuation. Effects and physical
channels absent from the formal event model are not claimed to be covered.
\end{abstract}

\section{Introduction}

Security mechanisms answer different questions.  A capability can answer
whether an executor possesses authority over an object.  An information-flow
policy can answer whether information at one security class may flow to
another.  A relational authorization rule can restrict which kind of executor
may become the next participant.  A provenance rule can restrict which
artifact, builder, or production history may be used.  These are not merely
different encodings of one scalar security property.  They are different
security state spaces.

The central proposal of this paper is to stop requiring one such state space to
subsume the others.  Instead, let each security system contribute one
coordinate to the state of a concrete execution occurrence, and let causal
continuity transport the product state.

Informally, if $S_1,\ldots,S_m$ are security-state systems, an occurrence at
step $i$ carries
\begin{equation}
  \state_i = (\coord_i^{(1)},\ldots,\coord_i^{(m)})
  \in \prod_{j=1}^{m} X_j,
\end{equation}
where $X_j$ is the state space of $S_j$.  Every system supplies its own
``no-more-permissive-than'' relation $\preceq_j$.  Ordinary causal continuation
requires
\begin{equation}
  \PoR(s_i,s_{i+1})
  \quad\land\quad
  \forall j,\;
  \coord_{i+1}^{(j)} \preceq_j \coord_i^{(j)}.
  \label{eq:core-transition}
\end{equation}
The first conjunct binds the states to one causal execution; the second prevents
any security coordinate from expanding silently while that execution
continues. At the event level, this is the lift from PIC's
$O\times R\times L$ to $\mathcal{B}\times L$: the causal axis is preserved,
while the behavior judged at an occurrence may now carry facets for every
security system in the product (Section~\ref{sec:leap}).

This generalizes the continuity pattern already present in
PIC.  The first PIC model carries an authority context $C_i\subseteq O\times R$
and requires $C_{i+1}\subseteq C_i$ along a PoR-linked execution
\cite{gallo_poc2026}.  Its relational extension already replaces the scalar
authority coordinate by the product state $(C_i,[\rho_i]_{\mathrm{sem}})$ and
orders that product componentwise by set inclusion and logical implication
\cite{gallo_relational2026}.  \Model takes the next step: the product itself is
the general object, while authority is one possible coordinate.

The contribution is not the observation that security systems can be combined.
The contribution is a minimal mathematical condition under which their states
can be \emph{continued together without collapsing their distinct semantics}:
causal linkage plus componentwise non-expansion.


\paragraph{Positioning.}
PIC was developed independently, starting from the propagation of authority
across distributed services and AI agents, and not from the language-based
information-flow or object-capability literature. Its primitives, execution
occurrences, an explicit lineage axis, receiver-verifiable PoR, non-expansive
carried authority, and a successor that need not exist when its predecessor
acts, were introduced for that setting \cite{gallo_poc2026,gallo_relational2026}.
The present paper takes those primitives as the foundation and asks a further
question: can the same continuity structure serve as a general theory, inside
which established security models can be placed and their results reread?
The readings of Bell--LaPadula, Denning, Boebert, Miller--Yee--Shapiro and
Rajani--Garg--Rezk that follow are therefore interpretations made from the
PIC/\Model frame. They do not claim that those authors used this frame, and
they do not make PIC a derivation of their work.

The frame also covers settings outside the scope of those works: authority
continuity across services and trust domains rather than within one program
semantics, continuity that the receiving boundary can verify, successors that
are unknown when a continuation is emitted, and a product of heterogeneous
security coordinates carried along one causal history.

A second objective is to close a boundary left explicit in the earlier PIC
series: causal attribution of an effect beyond an immediately witnessed
handoff. Within \Model, this requires an influence relation on execution
occurrences that represents every in-scope explicit and implicit dependency
and is contained in the causal order. Rajani, Garg, and Rezk (RGR) provide a
natural language-level instance of such a relation. Their full provenance
semantics (FPs) tracks all value influence represented by their calculus,
including implicit influence through branch and loop control, and proves
CDA-freedom for the entire calculus without the program and heap restrictions
required by their capability or explicit-only provenance results
\cite{rajani_garg_rezk2016}. We read that influence as one realization of the
\Model influence relation and place it on the same causal axis as PoR. The
purpose is not to replace RGR's theorem, but to show that, seen from PIC,
full provenance and authority continuity are two views of one event structure
and can therefore be composed into a single causal proof.

\paragraph{Contributions.}
This paper develops ten results.
\begin{enumerate}[leftmargin=*]
  \item An explicit lift of PIC's event space $O\times R\times L$ to a
        general security-event space $\mathcal{B}\times L$.
  \item A single underlying causal event order on which authority, IFC,
        provenance, relational state, and other security views are attached to
        the same execution occurrences.
  \item A general ordered security-state coordinate and its product
        construction, with a joint product-admission semantics.
  \item A Product Continuity theorem showing that every ordinary descendant
        is origin-bounded in every coordinate and jointly in admitted behavior.
  \item A per-coordinate continuity construction showing that PIC authority
        continuity and IFC/provenance continuity share PoR as the same causal
        basis; product continuity is their conjunction, not a second clock.
  \item An application of the principal down-set embedding and a
        projection/factorization theorem showing what is lost when a decision
        omits a security coordinate or causal history.
  \item A translated-continuity theorem for heterogeneous security-state
        spaces and a segment-continuity theorem for explicit authorized
        relaxations.
  \item A PIC-based reading of RGR's capability, explicit provenance, and
        full provenance results, faithful to their original statements,
        including their unconditional Theorem~6 for CDA-freedom with
        endorsement.
  \item A conservative event refinement showing that RGR's full-provenance
        influence is one instance of a \Model influence relation, causally
        ordered on the same occurrences as PoR; the derivation of causal
        attribution inside a closed event model depends only on the closure
        properties (E1)--(E3), not on RGR's calculus.
  \item A request-bounded effect theorem and a model-level exclusion of
        lineage-form confused-deputy states along ordinary continuation, with
        endorsement, composition, or authority relaxation represented as
        explicit authorized causal events.
\end{enumerate}

\section{Rajani--Garg--Rezk: Full Provenance and CDA Freedom}
\label{sec:rgr}

We summarize RGR's results because they give an independent, language-based
treatment of the confused deputy that can be read inside the PIC/\Model frame.
The summary is descriptive and follows their statements; the continuity and
attribution results of this paper are stated over the \Model event structure
and do not depend on their calculus.

Rajani, Garg, and Rezk study access control, object-capability semantics, and
confused-deputy attacks in a small imperative \emph{region calculus}
\cite{rajani_garg_rezk2016}. Their formal language contains principals,
mutable references, read and write capabilities $R_r$ and $W_r$, conditionals,
loops, assignments, sequencing, and commands executed inside regions. A
principal lattice supplies the write-integrity policy through an ownership map
$O$ from references to principals.

Their first comparison equips the same calculus with two operational semantics.
Access-control semantics checks authority when a reference is written.
Capability semantics implements Miller et al.'s Property~A, ``no designation
without authority,'' by preventing a region from obtaining a write capability
that it is not authorized to wield. RGR prove that the access-control semantics
is strictly more permissive than this capability semantics, while both satisfy
the write-integrity property they formulate.

The paper's second contribution is an extensional characterization of
confused-deputy freedom. An authority context $E_{\rho_A}$ contains a hole in
an adversary region $\rho_A$, and an attacker-interest set $AIS$ identifies the
references of concern. Informally, a reference is CDA-safe if either the
adversary cannot control its final value or the adversary could have written
that final value directly. Endorsed regions are then excluded from the attack
condition because an endorsed principal explicitly chooses to act on lower
integrity influence. RGR denote the resulting property by
$\CDAF(E_{\rho_A},H,\to_{red})$.

This definition lets RGR separate explicit from implicit designation. Their
capability semantics blocks the classical case in which the attacker supplies
the privileged reference itself, but it does not block all cases in which the
deputy already possesses the privileged reference and the attacker instead
controls the value written or the control path that determines the write.
Their Theorem~4 therefore proves CDA-freedom for capability semantics only
under strong restrictions on programs and the initial heap
\cite{rajani_garg_rezk2016}.

RGR then add provenance tracking to the access-control baseline. Their
explicit-only provenance semantics (EPs) labels each computed value by a
principal that lower-bounds the principals whose references have explicitly
influenced the value. EPs deliberately ignores implicit control-flow
influence, and their Theorem~5 still needs restrictions, although weaker ones
than capability semantics.

Their full provenance semantics (FPs) removes those remaining restrictions.
Every computed value still carries a principal-valued provenance label. In
addition, the semantic state carries a stack $PC$ of program-counter labels.
On entry to an \texttt{if} or \texttt{while} body, the new program-counter
label is the meet of the current $pc$ and the label of the condition; on exit,
the old $pc$ is restored. The FPs assignment rule checks the provenance of the
target, the value, and the current $pc$ against the owner of the reference being
written. Thus explicit data influence and implicit control influence are both
part of the enforcement semantics \cite{rajani_garg_rezk2016}.

The key result is their Theorem~6:
\begin{equation}
  \CDAF(E_{\rho_A},H,\to_{\FP})
  \quad\text{holds unconditionally.}
  \label{eq:rgr-thm6}
\end{equation}
``Unconditionally'' here is relative to the complete RGR region calculus and
FPs semantics: unlike their capability and EP results, no additional program or
heap restrictions are assumed. The proof route is also explicit. FPs
establishes a strong non-interference/information-flow-integrity property, and
RGR's Theorem~7 proves that non-interference with endorsement implies
CDA-freedom with endorsement \cite{rajani_garg_rezk2016}.

Two boundaries matter for the present paper. First, RGR's FPs labels are
\emph{principal-valued abstractions of influence}; their paper does not attach
event identifiers to labels. Second, RGR do not define PoR, PoC, a distributed
lineage protocol, or a Lamport-style event graph. The event construction below
is therefore a new conservative refinement of the influence already represented
by FPs, not a claim about terminology or machinery present in RGR.

\section{Historical Motivation: Different Security Coordinates}

Bell and LaPadula developed a mathematical model of secure computer states and
rules intended to preserve security under state transitions
\cite{bell_lapadula1973,bell_lapadula1973model}.  Denning subsequently formulated secure information
flow using a lattice of security classes, giving a general mathematical view of
systems that restrict information flow \cite{denning1976}.

Capability systems start from a different primitive. Authority is represented
by protected designators that permit operations on objects. In 1984, Boebert
showed that an \emph{unmodified} capability machine, in the absence of
additional mechanisms, cannot enforce the Bell--LaPadula $*$-property in his
model \cite{boebert1984}. His concrete attack lets a low-clearance program place
a read/write capability for a low object into that low object; a later
high-clearance program may read the low object, obtain the capability, and then
combine read access to high data with write access to the low object.  Boebert's
concluding explanation is that, in the pure capability machine he defines, the
right to exercise access carries the ability to propagate that access.

Miller, Yee, and Shapiro revisited this construction directly in
\emph{Capability Myths Demolished} \cite{miller_yee_shapiro2003}.  They agree
that the particular construction described by Boebert fails to enforce the
$*$-property, but argue that the attack assumes capabilities can be transmitted
wherever data can be transmitted.  In partitioned or type-enforced capability
systems, capability transfer and data transfer are distinguishable operations.
Their two-level construction independently controls those operations and is
presented as a counterexample to the broader claim that capability systems as a
class cannot enforce the $*$-property.  The same paper also stresses that a
static access matrix is not a complete description of a security mechanism:
the rules by which access relationships evolve over time materially affect the
security properties of the resulting system.

Miller and Shapiro develop a complementary argument in \emph{Paradigm
Regained} \cite{miller_shapiro2003}.  They distinguish permission from
authority and explicitly make program behavior on permitted causal pathways
part of authority analysis.  Their object-capability model carries overt
causation by messages over permitted paths.  Their examples use Redell's
Caretaker pattern for revocation, factories for overt confinement, and a data
diode for the $*$-properties.  They also
show how multiple independently motivated abstractions can coexist over the
same graph under ``mutually suspicious composition.''

These papers therefore already contain dynamic, behavioral, and causal
reasoning; \Model does not claim otherwise.  The distinction made here is
representational.  In those analyses, causal pathways and abstraction behavior
are derived from the reference graph and program semantics.  \Model instead
makes causal lineage an explicit axis of the security event and transports an
ordered security state along that axis.  Nor does \Model assume that a
capability implementation can contribute only an authority coordinate: a
capability design may itself mediate information flow, transfer type, request
context, or another security distinction.  The general claim is narrower: if an
enforcement decision is invariant under a distinction that changes the target
security verdict, that decision cannot enforce the target property.  A typed
transfer rule or access abstraction may expose exactly the distinction that a
coarser model omitted.

This paper does \emph{not} claim that Boebert, Bell--LaPadula, Denning, Miller,
or Shapiro formulated their work in this product language.  The product view is
a new interpretation intended to explain why apparently conflicting results
can all be correct under different observable state spaces.

\section{Preliminaries}

\subsection{Execution occurrences and the common causal order}

Let $S$ be the set of concrete execution occurrences of one modeled run. The
machine may be sequential, concurrent, or distributed. Let
\begin{equation}
  \to_c\;\subseteq S\times S
\end{equation}
be an acyclic immediate causal-predecessor relation, and let $\hb$ denote its
transitive closure. In a sequential reduction trace, $\to_c$ is a chain. In a
concurrent or distributed execution it is generally a DAG, in the sense of a
causal happened-before style order rather than a requirement of one global
clock \cite{lamport1978}. This causal order is the underlying order of the run;
security mechanisms do not receive independent physical clocks merely because
they use different state spaces. We require the semantic causal order to be
\emph{dependency adequate}: whenever an occurrence consumes a value, message,
or state produced by an earlier occurrence, or its reachability is
control-dependent on a condition evaluated at an earlier occurrence, the
producer/condition is an ancestor of the consumer/controlled occurrence under
$\hb$. This is a local semantic condition on causal edges; it does not assume
the request-to-effect Causal Attribution theorem proved later.

PoR is the receiver-verifiable immediate-continuation relation used by PIC and
\Model. At the semantic level considered here, sound PoR satisfies
\begin{equation}
  \PoR(s,s')\Longrightarrow s\to_c s'.
  \label{eq:por-sound}
\end{equation}
For a closed semantic execution, every immediate causal edge in the modeled
security-relevant event graph is required to be represented by a PoR-witnessable
transition. A concrete realization may establish this relation using signatures,
attestations, observed handoffs, trusted-runtime evidence, or an
information-flow/provenance witness. Evidence and the relation evidenced are
kept distinct.

A finite \emph{lineage path} is
\begin{equation}
  \pi=\langle s_0,s_1,\ldots,s_n\rangle
\end{equation}
with $\PoR(s_i,s_{i+1})$ at every adjacent step. Let $L$ denote the set of
such finite causal paths. In a DAG, several lineage paths may share an origin
or endpoint; the underlying occurrence $s\in S$ is unique even when it belongs
to several paths. The path notation is retained because the basic continuity
proof is transitivity along a causal path.

\subsection{One occurrence, several security views}

An occurrence $s_i$ may simultaneously carry authority, information-flow,
provenance, identity, relational, geographic, or other security observations.
These coordinates may inhabit unrelated mathematical spaces, but they are
observations of the same event. For example,
\begin{equation}
  \omega_i=(\ell_i,\iota_i,g_i,\state_i),
\end{equation}
where $\ell_i$ is causal position or lineage information, $\iota_i$ is an
executor or identity observation, $g_i$ is a geographic or network location,
and $\state_i$ is the security state carried by the occurrence.

Lineage grows by causal extension rather than attenuation, while the current
identity or location may vary from step to step. If geography is a security
restriction, the monotone object is instead an admissible set or predicate, e.g.
\begin{equation}
  g_i\in G_i^{\mathrm{allow}},
  \qquad
  G_{i+1}^{\mathrm{allow}}\subseteq G_i^{\mathrm{allow}}.
\end{equation}
This distinction prevents the causal coordinates of an occurrence from being
confused with the policy state that constrains them.

\subsection{Per-coordinate continuity over the same PoR}
\label{sec:pocj}

For a carried coordinate $j$, define its path-continuity property by
\begin{equation}
 \PoC_j(\pi)
 \;\Longleftrightarrow\;
 \bigwedge_{i=0}^{n-1}
 \left(
   \PoR(s_i,s_{i+1})\land
   \coord_{i+1}^{(j)}\preceq_j\coord_i^{(j)}
 \right).
 \label{eq:pocj}
\end{equation}
The authority-only instance is the original PIC/PoC condition
$C_{i+1}\subseteq C_i$. An IFC coordinate has its own refinement order, while
full provenance contributes an influence relation on the same occurrences
(Section~\ref{sec:fp-axis}). Product continuity is therefore not a product of
independent times. It is the conjunction of several continuity disciplines all
sharing the same PoR edges:
\begin{equation}
  \PoC_{\mathrm{prod}}(\pi)
  \Longleftrightarrow
  \bigwedge_{j\in J}\PoC_j(\pi).
  \label{eq:pocprod}
\end{equation}
This formalizes the distinction between the common causal substrate and the
coordinate-specific security invariants transported over it.

\subsection{From \texorpdfstring{$O\times R\times L$}{O x R x L} to \texorpdfstring{$\mathcal{B}\times L$}{B x L}}
\label{sec:leap}

In PIC an authority event is an occurrence of a privilege within a
causal lineage,
\begin{equation}
  e=(o,r,\ell)\in O\times R\times L,
\end{equation}
and the state carried by an occurrence is an authority context
$C\subseteq O\times R$ \cite{gallo_poc2026}. The privilege coordinates
$O\times R$ say \emph{what} is exercised; the lineage coordinate $L$ says
\emph{in which causal history} it is exercised. \Model keeps the lineage
axis unchanged and replaces the single privilege universe by the product
of the behavior universes of all carried security systems.

Anticipating Definition~\ref{def:coordinate}, let each carried security
system contribute a behavior universe $\mathcal{B}_j$.

\begin{definition}[Security event]\label{def:event}
Let $\mathcal{B}=\prod_{j\in J}\mathcal{B}_j$. A security event is
\begin{equation}
  e=(b,\ell)\in\mathcal{B}\times L,
\end{equation}
where $\ell$ is the causal history of the occurrence and
$b=(b^{(j)})_{j\in J}$ collects, for each coordinate, the facet of the
occurrence that the coordinate judges: the privilege it exercises, the
flow it performs, the executor transition it realizes, or the location at
which it occurs.
\end{definition}

For the single authority coordinate, $\mathcal{B}_{\mathrm{auth}}=O\times
R$ and $\mathcal{B}\times L=O\times R\times L$: PIC's event space is the
one-coordinate instance. The same substitution acts on states, where
$\powerset(O\times R)$ becomes $X=\prod_j X_j$. The construction of PIC
is otherwise transported unchanged: states are carried along PoR-linked
occurrences, and ordinary continuation may only refine them.

\begin{remark}
If a coordinate imposes nothing on a given occurrence, its facet may be
a distinguished element $\bot_j\in\mathcal{B}_j$ admitted by every state
of that coordinate.
\end{remark}

\section{Ordered Security-State Systems}

\begin{definition}[Security-state coordinate]\label{def:coordinate}
A security-state coordinate is a tuple
\begin{equation}
  \mathfrak{S}_j=(X_j,\preceq_j,\mathcal{B}_j,\sem{\cdot}_j),
\end{equation}
where:
\begin{itemize}[leftmargin=*]
  \item $X_j$ is a set of security states;
  \item $\preceq_j$ is a partial order interpreted as ``no more permissive
        than'';
  \item $\mathcal{B}_j$ is a universe of behaviors, transitions, principals,
        flows, or other semantically admitted possibilities; and
  \item $\sem{x}_j\subseteq\mathcal{B}_j$ is the semantic extension admitted
        by state $x\in X_j$.
\end{itemize}
The coordinate is \emph{semantically monotone} when
\begin{equation}
  x'\preceq_j x
  \quad\Longrightarrow\quad
  \sem{x'}_j\subseteq\sem{x}_j.
  \label{eq:semantic-monotonicity}
\end{equation}
\end{definition}

The definition deliberately does not require every security system to be born
as a powerset.  A lattice label, a logical predicate, a capability set, a
provenance constraint, or a history-dependent state can all be used if they
admit a sound refinement order.

\begin{remark}[Coordinates are semantic dimensions]
A coordinate is not required to correspond one-to-one with an implementation
mechanism.  One mechanism may realize several security coordinates, as when a
type-enforced capability runtime distinguishes authority from data transfer and
capability transfer; conversely, several mechanisms may jointly realize one
coordinate.  \Model reasons over the security distinctions exposed to the
admission decision, not over product names or implementation categories.
\end{remark}

\begin{remark}[Refinement is not a complete update semantics]
The order $\preceq_j$ is the condition needed for the continuity theorem, not
a claim that every refinement step is reachable or permitted by the concrete
security system.  A coordinate may impose a stricter local transition relation
$U_j\subseteq X_j\times X_j$ satisfying
\begin{equation}
  U_j(x,x')\Longrightarrow x'\preceq_j x.
\end{equation}
All ordinary-continuation results below remain valid when transitions must
satisfy both $U_j$ and the refinement order.  This leaves room for the richer
update rules that distinguish dynamic security mechanisms even when they share
the same static state representation.
\end{remark}

\begin{definition}[Product security state]
For a finite family $J$ of security-state coordinates, define
\begin{equation}
  X \;=\; \prod_{j\in J} X_j,
  \qquad
  \state=(\coord^{(j)})_{j\in J}.
\end{equation}
The product refinement order is
\begin{equation}
  \state'\preceq\state
  \quad\Longleftrightarrow\quad
  \forall j\in J,\;
  \coord'^{(j)}\preceq_j\coord^{(j)}.
  \label{eq:product-order}
\end{equation}
\end{definition}

For a concrete execution, let
\begin{equation}
  \st:S\to X
\end{equation}
assign exactly one product security state to each execution occurrence, and
write $\state_i:=\st(s_i)$. If $\ell\in L$ is a finite causal history with
endpoint $s_n$, define
\begin{equation}
  \state(\ell):=\st(s_n).
  \label{eq:lineage-endpoint-state}
\end{equation}
Thus the state used by an event is a well-defined function of the endpoint
occurrence of its represented history.

\begin{proposition}[Product order]
If every $(X_j,\preceq_j)$ is a partial order, then
$(X,\preceq)$ defined by \eqref{eq:product-order} is a partial order.
\end{proposition}

\begin{proof}
Reflexivity, antisymmetry, and transitivity hold componentwise.
\end{proof}

\begin{remark}[Coordinates need not be independent]
The Cartesian product is an ambient state space, not an assumption that every
combination of coordinate states is meaningful. If cross-coordinate constraints
permit only states in $Y\subseteq X$, the induced order on $Y$ is still a
partial order and Product Continuity applies unchanged to lineages whose states
remain in $Y$. Likewise, if only certain combinations of behavior facets are
meaningful, one may restrict the event universe to
$\mathcal{B}^{*}\subseteq\prod_j\mathcal{B}_j$ and interpret product admission
on $\mathcal{B}^{*}\cap\prod_j\sem{\coord^{(j)}}_j$.
\end{remark}

\section{Causally Ordered Security-State Continuity}

\begin{definition}[Valid \Model transition]
A transition
\begin{equation}
  (s_i,\state_i)\longrightarrow(s_{i+1},\state_{i+1})
\end{equation}
is \Model-valid when
\begin{equation}
  \PoR(s_i,s_{i+1})
  \quad\land\quad
  \state_{i+1}\preceq\state_i.
  \label{eq:valid-cosst}
\end{equation}
Equivalently, every security coordinate refines or preserves its predecessor
state while the execution advances along a verified causal edge.
\end{definition}

\begin{definition}[Product security-state continuity]
A lineage
\begin{equation}
  \Pi=\langle(s_0,\state_0),\ldots,(s_n,\state_n)\rangle
\end{equation}
has product security-state continuity when every adjacent transition is
\Model-valid.
\end{definition}

\begin{theorem}[Product Continuity]
If $\Pi$ has product security-state continuity, then
\begin{equation}
  \state_n\preceq\state_{n-1}\preceq\cdots\preceq\state_0.
\end{equation}
Hence, for every coordinate $j$,
\begin{equation}
  \coord_n^{(j)}\preceq_j\coord_0^{(j)}.
  \label{eq:component-origin-bound}
\end{equation}
If the coordinate is semantically monotone, then
\begin{equation}
  \sem{\coord_n^{(j)}}_j
  \subseteq
  \sem{\coord_0^{(j)}}_j.
\end{equation}
\end{theorem}

\begin{proof}
Every valid transition gives
$\state_{i+1}\preceq\state_i$.  Transitivity of the product order gives
$\state_n\preceq\state_0$.  Projection to coordinate $j$ gives
\eqref{eq:component-origin-bound}; semantic inclusion follows from
\eqref{eq:semantic-monotonicity}.
\end{proof}

The theorem states the general continuity invariant for lineages containing
only ordinary continuations: causal transport cannot introduce a possibility
forbidden by the origin state of any carried security system.


\begin{theorem}[Pathwise Product Continuity on a causal DAG]
\label{thm:dag-continuity}
Let $q,e\in S$ with $q\hb e$. If every edge of a causal path from $q$ to $e$
is an ordinary \Model transition, then
\begin{equation}
  \st(e)\preceq \st(q).
\end{equation}
If every immediate predecessor edge of a multi-parent ordinary join is required
to satisfy the same componentwise non-expansion rule, the bound holds along
every parent path.
\end{theorem}

\begin{proof}
Choose a causal path
$q=s_0\to_c s_1\to_c\cdots\to_c s_k=e$. Each ordinary edge gives
$\st(s_{i+1})\preceq\st(s_i)$. Transitivity yields
$\st(e)\preceq\st(q)$. The multi-parent statement follows because the same
argument applies independently to every immediate-predecessor path.
\end{proof}

\begin{definition}[Product admission]\label{def:product-admission}
The product state $\state=(\coord^{(j)})_{j\in J}$ admits
$b\in\mathcal{B}$ iff $b^{(j)}\in\sem{\coord^{(j)}}_j$ for every $j$.
Write
\begin{equation}
  \sem{\state}:=\prod_{j\in J}\sem{\coord^{(j)}}_j\subseteq\mathcal{B}.
\end{equation}
\end{definition}

\begin{proposition}[Joint semantic monotonicity]\label{prop:joint}
If every coordinate is semantically monotone, then
$\state'\preceq\state$ implies $\sem{\state'}\subseteq\sem{\state}$.
Consequently, under product security-state continuity,
$\sem{\state_n}\subseteq\sem{\state_0}$.
\end{proposition}

\begin{proof}
$\state'\preceq\state$ gives $\coord'^{(j)}\preceq_j\coord^{(j)}$ for
every $j$, hence $\sem{\coord'^{(j)}}_j\subseteq\sem{\coord^{(j)}}_j$ by
semantic monotonicity, and a product of inclusions is an inclusion. The
second claim follows from Product Continuity.
\end{proof}

A later occurrence therefore admits only behaviors that the origin state
admits in every coordinate jointly, not merely coordinate by coordinate.

\section{Why Refinement Can Be Read as Set Inclusion}

The previous theorem is intentionally stated over partial orders rather than
subsets.  Nevertheless, the subset intuition can be recovered for every
partial order.

\begin{definition}[Principal lower-set representation]
For a poset $(X,\preceq)$ define
\begin{equation}
  \downarrow x
  =\{y\in X\mid y\preceq x\}.
\end{equation}
\end{definition}

\begin{theorem}[Principal down-set embedding]
For every poset $(X,\preceq)$, the standard principal down-set map
\begin{equation}
  \eta:X\to\powerset(X),
  \qquad
  \eta(x)=\downarrow x
\end{equation}
is an order embedding:
\begin{equation}
  x'\preceq x
  \quad\Longleftrightarrow\quad
  \downarrow x'\subseteq\downarrow x.
  \label{eq:downset-embedding}
\end{equation}
\end{theorem}

\begin{proof}
If $x'\preceq x$ and $y\preceq x'$, transitivity gives $y\preceq x$, so
$\downarrow x'\subseteq\downarrow x$.  Conversely,
$\downarrow x'\subseteq\downarrow x$ implies $x'\in\downarrow x$ because
$x'\in\downarrow x'$ by reflexivity.  Hence $x'\preceq x$.
\end{proof}

This standard principal down-set embedding formalizes the subset intuition
behind the present work. A security system need not literally store a set, but
its refinement order can be represented faithfully as inclusion among sets of
states no more permissive than the current one. This inclusion must not be
confused with semantic admission: $\downarrow x\subseteq X$ is a set of
\emph{states}, whereas $\sem{x}\subseteq\mathcal{B}$ is a set of
\emph{behaviors}. Semantic monotonicity connects refinement to behavior
inclusion, but the two constructions are distinct. The familiar PIC rule
$C_{i+1}\subseteq C_i$ is therefore one concrete instance of a more general
ordered refinement pattern.

\section{Projection, Quotients, and Coordinate Blindness}

The second PIC paper interprets lineage-forgetting as a quotient and shows that
lineage-invariant decisions factor through the projection that forgets causal
history \cite{gallo_relational2026}.  The same construction applies to arbitrary
products of security-state coordinates.

Let
\begin{equation}
  X=\prod_{j\in J}X_j
\end{equation}
and let $K\subseteq J$.  Define the coordinate projection
\begin{equation}
  p_K:X\to X_K:=\prod_{j\in K}X_j.
\end{equation}

\begin{definition}[$K$-invariant decision]
A decision rule $A:X\to\{0,1\}$ is $K$-invariant when
\begin{equation}
  p_K(x)=p_K(y)
  \quad\Longrightarrow\quad
  A(x)=A(y).
\end{equation}
\end{definition}

\begin{theorem}[Projection Factorization]
A decision rule $A:X\to\{0,1\}$ is $K$-invariant if and only if there exists a
unique rule
\begin{equation}
  \bar A:p_K(X)\to\{0,1\}
\end{equation}
such that
\begin{equation}
  A=\bar A\circ p_K.
  \label{eq:factorization}
\end{equation}
\end{theorem}

\begin{proof}
If $A=\bar A\circ p_K$, equal projections yield equal decisions.  Conversely,
define $\bar A(z)=A(x)$ for any $x$ with $p_K(x)=z$.  $K$-invariance makes the
definition independent of the representative.  Uniqueness follows because
$p_K$ is surjective onto $p_K(X)$.
\end{proof}

\begin{corollary}[Omitted-coordinate impossibility]
Suppose a target security property $P:X\to\{0,1\}$ distinguishes two states
$x,y\in X$ for which $p_K(x)=p_K(y)$.  No $K$-invariant decision rule can equal
$P$ on all of $X$.
\end{corollary}

\begin{proof}
Every $K$-invariant rule must give the same result on $x$ and $y$, while $P$
requires different results.
\end{proof}

This is the general form of the earlier lineage-forgetting observation: a
security decision cannot enforce a distinction that its observation function
erases.

\begin{corollary}[Behavioral-state blindness]\label{cor:behavioral-blindness}
Let $X=X_{\mathrm{prot}}\times X_{\mathrm{beh}}$ separate a primitive
protection-state coordinate from a coordinate describing security-relevant
behavior or an access abstraction.  If a target property $P:X\to\{0,1\}$
satisfies
\begin{equation}
  P(p,h_1)\ne P(p,h_2)
\end{equation}
for some fixed $p\in X_{\mathrm{prot}}$, then no decision rule that factors
through the projection $(p,h)\mapsto p$ can implement $P$ on all of $X$.
\end{corollary}

\begin{proof}
This is Omitted-coordinate Impossibility with
$K=\{\mathrm{prot}\}$.
\end{proof}

This corollary gives an order-independent statement of the blind spot discussed
by Miller and Shapiro: if two systems have the same protection graph but the
behavior of a security-enforcing abstraction changes the authority or flow
verdict, an analysis that observes only the protection graph is necessarily too
coarse for that distinction \cite{miller_shapiro2003}.  It does not imply that
protection-state analysis is unsound when it succeeds; it states only that such
a projection cannot reproduce a behavior-sensitive property in general.

Using the endpoint-state assignment of \eqref{eq:lineage-endpoint-state}, the
\Model admission rule for an event $(b,\ell)$ is
\begin{equation}
  A_{\Model}(b,\ell)=1
  \quad\Longleftrightarrow\quad
  b\in\sem{\state(\ell)}.
\end{equation}
Let $q_L:\mathcal{B}\times L\to\mathcal{B}$, $(b,\ell)\mapsto b$, forget
the lineage.

\begin{corollary}[Lineage-forgetting blindness in \Model]
\label{cor:coss-lineage}
Let $b\in\sem{\state(\ell)}\setminus\sem{\state(\ell')}$. No decision rule
that factors through $q_L$ equals $A_{\Model}$ on both $(b,\ell)$ and
$(b,\ell')$.
\end{corollary}

\begin{proof}
$q_L(b,\ell)=q_L(b,\ell')$, so every rule factoring through $q_L$ gives
both events the same verdict, while $A_{\Model}(b,\ell)=1$ and
$A_{\Model}(b,\ell')=0$. This is the same fiber argument as Projection
Factorization, now applied to the projection that omits $L$ from the event
space.
\end{proof}

For the authority coordinate alone this is the lineage-invariance result
of PIC \cite{gallo_poc2026}. In \Model it holds for every carried security
system at once: the same behavior can be admissible in one causal history
and inadmissible in another, in authority, in information flow, in
relational admissibility or in any other coordinate, and a decision that
forgets the history cannot tell the two apart.

\subsection{A three-coordinate example}
Consider a three-step agent lineage
\begin{equation}
  s_0\to s_1\to s_2
\end{equation}
carrying authority, information-flow, and geographic coordinates
\begin{equation}
  \state_i=(C_i,F_i,G_i),
  \qquad
  \state_2\preceq\state_1\preceq\state_0.
\end{equation}
At $s_2$, suppose $C_2$ still permits the operation
$(\mathrm{send},\mathrm{reportAPI})$ and $G_2$ admits the agent's current EU
execution region, but $F_2$ excludes the flow facet $f_{\mathrm{ext}}$
representing transmission of the current data to an external sink. Let $b$ be
the event facet tuple whose authority component is that send operation, whose
flow component is $f_{\mathrm{ext}}$, and whose geographic component is the
current EU region. Authority and geography admit their respective facets, but
IFC does not, so
\begin{equation}
  b\notin\sem{\state_2}
\end{equation}
and the product rejects the event. An authority-only decision would admit the
event even though the complete product state rejects it: the authority
coordinate itself remains valid, while the IFC coordinate does not.

Now consider a different causal history $\ell'$ carrying the same authority
and geographic states but an IFC state $F'$ that admits $f_{\mathrm{ext}}$,
for example after a separately authorized declassification. If $\ell$ denotes
the first history ending at $s_2$, the same event facet tuple $b$ may satisfy
\begin{equation}
  b\in\sem{\state(\ell')}
  \qquad\text{while}\qquad
  b\notin\sem{\state(\ell)}.
\end{equation}
The example exhibits both claims of the theory: all relevant security
coordinates constrain one causal occurrence jointly, and a decision that
forgets the lineage cannot distinguish histories whose current admissibility
differs.

\section{Classical Security Models as Coordinates}

\subsection{Authority and capabilities}

Let $O$ be operations and $R$ resources.  Define
\begin{equation}
  X_{\mathrm{auth}}=\powerset(O\times R),
  \qquad
  C'\preceq_{\mathrm{auth}}C
  \Longleftrightarrow C'\subseteq C.
\end{equation}
This coordinate says which authority may participate in the current causal
execution.  It does not identify the executor that happens to possess other
ambient authority.

With this coordinate alone, \Model-validity reduces to
\begin{equation}
  \PoR(s_i,s_{i+1})
  \quad\land\quad
  C_{i+1}\subseteq C_i,
\end{equation}
which is precisely the abstract continuity rule of PIC
\cite{gallo_poc2026}.

This specialization is compatible with the stronger object-capability results
in the literature.  \emph{Capability Myths Demolished} identifies absence of
ambient authority, inseparability of designation and authority, and
access-controlled delegation channels as important properties of
object-capability systems; it also emphasizes that authority conveyed in the
context of a request helps avoid the classical confused-deputy pattern
\cite{miller_yee_shapiro2003}.  \Model does not treat those protections as
missing.  A capability system may realize several security distinctions at
once; the lineage coordinate matters only for a target property whose verdict
also varies with causal execution after the capability-level distinctions have
been fixed.

\subsection{Information-flow admissibility}

Let $D$ be a domain of information sources and sinks, security classes, or
other flow endpoints.  One possible semantic IFC coordinate is the set of
currently admitted flows
\begin{equation}
  F\subseteq D\times D.
\end{equation}
Define
\begin{equation}
  F'\preceq_{\mathrm{ifc}}F
  \Longleftrightarrow F'\subseteq F.
  \label{eq:ifc-subset}
\end{equation}
A Bell--LaPadula or lattice policy can induce such an admissible-flow relation;
this paper does not replace the internal semantics of those models.  It treats
their effective admissibility state as an independent coordinate whose ordinary
continuation must not become more permissive.

An alternative is to take lattice labels or complete IFC policy states as
$X_{\mathrm{ifc}}$ and define $\preceq_{\mathrm{ifc}}$ directly from their
semantic flow sets.  Equation \eqref{eq:ifc-subset} is therefore an example,
not a claim that all IFC implementations literally store $F$.


The policy state $F_i$ should be distinguished from the dynamic provenance of a
particular run. The former answers which flows are admissible; the latter
answers which earlier occurrences actually influenced a value or control
decision in this execution. RGR's FPs supplies a formal influence abstraction
through value labels and the program-counter stack. Section~\ref{sec:fp-axis}
lifts that influence to event resolution without changing the original FPs
accept/reject behavior.

\subsection{Typed data and capability transfer}

The distinction used by Miller, Yee, and Shapiro to revisit Boebert can itself
be represented as an ordered security coordinate without claiming that their
systems used this notation.  Let $\Lambda$ be a set of security domains and
let
\begin{equation}
  K=\{\mathsf{data},\mathsf{cap}\},
  \qquad
  \mathcal{B}_{\mathrm{tr}}=\Lambda\times\Lambda\times K
\end{equation}
classify a transfer by source domain, destination domain, and whether the
transferred value is ordinary data or a capability.  Define
\begin{align}
  X_{\mathrm{tr}}&=\powerset(\mathcal{B}_{\mathrm{tr}}),\\
  Q'\preceq_{\mathrm{tr}}Q
  &\Longleftrightarrow Q'\subseteq Q,
  \qquad
  \sem{Q}_{\mathrm{tr}}=Q.
\end{align}
The coordinate can therefore admit a data transfer between two levels while
rejecting a capability transfer over the same pair of levels.

This captures the specific distinction made in \emph{Capability Myths
Demolished}: Boebert's attack requires a high-clearance subject to obtain, by
reading a low object, a capability that confers write access to that low object;
the later type-enforced construction distinguishes capability movement from
data movement and blocks the cross-level capability transfer
\cite{miller_yee_shapiro2003}.  The coordinate $Q$ is only an illustrative
semantic encoding of that distinction, not a claim that an object-capability
implementation literally stores a set of typed flows.

\subsection{Relational admissibility}

The relational extension of PIC carries a predicate $\rho_i$ over
successor/predecessor profiles and orders predicates by implication
\cite{gallo_relational2026}.  After quotienting by semantic equivalence, define
\begin{equation}
  X_{\mathrm{rel}}
  =\mathrm{Pred}(\mathcal{T}_J)/{\equiv_{\mathrm{sem}}},
\end{equation}
with
\begin{equation}
  [\rho']\preceq_{\mathrm{rel}}[\rho]
  \quad\Longleftrightarrow\quad
  \rho'\Rightarrow\rho.
\end{equation}
Semantically,
\begin{equation}
  \rho'\Rightarrow\rho
  \quad\Longleftrightarrow\quad
  \sem{\rho'}\subseteq\sem{\rho}.
\end{equation}
Thus the existing PIC product
\begin{equation}
  \Gamma_i=(C_i,[\rho_i]_{\mathrm{sem}})
\end{equation}
is already a two-coordinate instance of \Model.

\subsection{Geographic, provenance, and product constraints}

A geographic restriction can be represented by an admissible region set
$G_i^{\mathrm{allow}}$ ordered by inclusion, while the observed current
location $g_i$ is checked for membership.  A software-product or provenance
coordinate can analogously represent the set of admissible artifact identities,
builders, attestations, or derivation histories.  These coordinates do not need
to share an ontology with authority or IFC; they need only expose a sound
refinement order.

This distinction is useful for invocation.  A product can participate in many
runtime invocations, just as one executor can realize many execution
occurrences.  Product identity is therefore not itself an invocation.  A future
DAG extension can treat invocation as a causal join between an authority
lineage and a product/provenance lineage while preserving the coordinates that
justify each parent.

\section{Full Provenance on the COSS Causal Axis}
\label{sec:fp-axis}

\Model needs an influence relation on execution occurrences: a relation that
records which earlier occurrences influence the data or control of later ones.
Any relation satisfying the closure properties (E1)--(E3) of
Definition~\ref{def:event-complete} serves. RGR's full provenance semantics is a
natural language-level instance, and this section shows how it fits.
FPs labels values and the current control context with principal-valued
integrity information, whereas \Model relates concrete execution occurrences and
transports security state along them. The connection is obtained by refining
an FPs reduction trace with ghost event dependencies and then erasing those
dependencies back to the original FPs semantics.

\subsection{Conservative event refinement of FPs}

Consider an RGR FPs execution
\begin{equation}
  \Theta_0\to_{\FP}\Theta_1\to_{\FP}\cdots\to_{\FP}\Theta_n,
  \label{eq:fp-trace}
\end{equation}
where each $\Theta_i$ contains the heap, the $PC$ stack, and the current
program/command state as in RGR. Associate a distinct occurrence $e_i\in S$
with each reduction step $\Theta_i\to_{\FP}\Theta_{i+1}$.

\begin{definition}[Event-refined full-provenance influence]
The relation $\Influence\subseteq S\times S$ is generated over a trace
\eqref{eq:fp-trace} by the following dependencies.
\begin{enumerate}[label=(F\arabic*),leftmargin=*]
  \item \emph{Explicit data/designation influence.} If a value, reference,
        target designation, or heap value produced or obtained at $e_i$ is a
        semantic input to the computation performed at a later event $e_j$,
        then $e_i\Influence e_j$.
  \item \emph{Implicit control influence.} If the reachability of the command
        executed at $e_j$ is control-dependent on a branch or loop condition
        evaluated at $e_i$, then $e_i\Influence e_j$.
\end{enumerate}
Its transitive closure is written $\Influence^{+}$.
\end{definition}

The first generator exposes at event resolution the explicit influence that
RGR propagate through expression provenance. The second exposes the implicit
influence that RGR propagate through the $PC$ stack. Event identifiers are
\emph{ghost annotations}: they do not occur in the premises of an RGR FPs
reduction rule.

\begin{proposition}[Conservative erasure]
Erasing event identifiers and $\Influence$ edges from an event-refined FPs run
yields exactly the underlying RGR FPs reduction trace, with the same
principal-valued labels, the same $PC$ stack evolution, and the same accepted
or blocked assignments.
\end{proposition}

\begin{proof}
The refinement adds only ghost metadata to the already existing reduction
steps. No original FPs premise reads that metadata. Erasure therefore commutes
with one reduction step; induction over the length of \eqref{eq:fp-trace} gives
the result.
\end{proof}

\begin{lemma}[Full-provenance influence is causally ordered]
\label{lem:fp-causal}
For every event-refined FPs run,
\begin{equation}
  e_i\Influence e_j\Longrightarrow e_i\hb e_j.
  \label{eq:fp-suborder}
\end{equation}
Consequently $\Influence^{+}\subseteq\hb$.
\end{lemma}

\begin{proof}
For explicit influence, the value/reference or state consumed at $e_j$ is
produced or obtained at $e_i$. By dependency adequacy of the semantic causal
order, producer/read events are ancestors of their consumers, hence
$e_i\hb e_j$. For implicit influence, RGR lower the current $pc$ only after
the branch or loop condition has been evaluated and execute the controlled
command while that lowered $pc$ is on the stack. The reachability of the
controlled event therefore depends on the condition event; dependency
adequacy again gives $e_i\hb e_j$. These are exactly the two generators of
$\Influence$. Transitivity of $\hb$ gives the closure statement.
\end{proof}

The reverse inclusion is neither needed nor generally true. Two commands may
be ordered by execution without the earlier one contributing data or control
influence to the later one. Full provenance therefore induces an influence
suborder on one common execution order; it does not create a second time.

\subsection{Event-complete provenance and derived causal attribution}

The sequential RGR calculus already closes all of its reads, writes, and
control constructs inside one operational semantics. For a general machine
execution we state the corresponding closure property directly.

\begin{definition}[Closed event-complete execution]
\label{def:event-complete}
A \Model execution is \emph{event-complete} for a stated threat model when:
\begin{enumerate}[label=(E\arabic*),leftmargin=*]
  \item every protected security-relevant effect is represented by an
        occurrence in $S$;
  \item every explicit data/designation dependency and every implicit control
        dependency in scope is represented by $\Influence^{+}$; and
  \item every generating influence edge is causally sound as in
        Lemma~\ref{lem:fp-causal}.
\end{enumerate}
For a request-origin occurrence $q$ and an effect occurrence $e$, define
\begin{equation}
  \CausedBy(e,q)
  \quad\Longleftrightarrow\quad
  q\Influence^{+}e.
  \label{eq:causedby-reach}
\end{equation}
\end{definition}

This definition is event based, not stack based. Synchronous calls, shared
state, a database, a queue, a timer, or an asynchronous callback do not require
a different causal principle. When the later event consumes state or a message
produced by the earlier event, the machine semantics contains the corresponding
dependency, for example
\begin{align}
 \mathrm{write}_x &\;\Influence\; \mathrm{read}_x,\\
 \mathrm{enqueue}_q &\;\Influence\; \mathrm{dequeue}_q,\\
 \mathrm{setTimer} &\;\Influence\; \mathrm{fireTimer},\\
 \mathrm{registerCallback} &\;\Influence\; \mathrm{invokeCallback}.
\end{align}
The wall-clock delay or disappearance of the original call stack does not erase
causal ancestry. What matters is whether the modeled later occurrence depends
on the earlier one.

\begin{theorem}[Causal Attribution from full-provenance reachability]
\label{thm:derived-ca}
In a closed event-complete execution,
\begin{equation}
  \CausedBy(e,q)\Longrightarrow q\hb e.
  \label{eq:ca-derived}
\end{equation}
Hence causal attribution is a theorem of the event semantics rather than an
independent assumption.
\end{theorem}

\begin{proof}
By Definition~\ref{def:event-complete}, $\CausedBy(e,q)$ is
$q\Influence^{+}e$. Lemma~\ref{lem:fp-causal} gives
$\Influence^{+}\subseteq\hb$. Therefore $q\hb e$.
\end{proof}

\subsection{Effect admission and authority continuity}

Let the authority coordinate be $C(s)\subseteq O\times R$. An ordinary causal
edge preserves PIC authority continuity:
\begin{equation}
  s\to_c s'\;\land\;\mathrm{Ord}(s,s')
  \Longrightarrow C(s')\subseteq C(s).
  \label{eq:authority-edge}
\end{equation}
At a multi-parent ordinary event the same rule is required for every immediate
predecessor; equivalently, the new authority is contained in the intersection
of the parent authority contexts. Any transition that legitimately broadens
authority is not an ordinary continuation: it is an explicitly authorized
origination, endorsement, composition, or relaxation event and starts a new
authority bound, as in Section~\ref{sec:relaxation}.

Let $Eff\subseteq S$ be the protected effect occurrences and let
\begin{equation}
  \Realize:Eff\to O\times R
\end{equation}
map an effect to the operation and actual resource it realizes. In the closed
model, admission of a protected effect is part of the semantics:
\begin{equation}
\begin{aligned}
  &e\in Eff\;\land\;\Realize(e)=(o,r)\;\land\;\Accepted(e)\\
  &\hspace{2.2em}\Longrightarrow (o,r)\in C(e).
\end{aligned}
\label{eq:effect-admission}
\end{equation}
Thus complete mediation is internalized in the formal semantics: protected
effects are precisely the events to which the admission rule applies.

\begin{lemma}[Request-origin authority bound]
\label{lem:req-bound}
Let $q\hb e$ and suppose the causal path from $q$ to $e$ contains only ordinary
authority-continuous edges. Then
\begin{equation}
  C(e)\subseteq C(q).
  \label{eq:req-bound}
\end{equation}
\end{lemma}

\begin{proof}
Choose a causal path
$q=s_0\to_c s_1\to_c\cdots\to_c s_k=e$. Equation
\eqref{eq:authority-edge} gives
$C(s_k)\subseteq C(s_{k-1})\subseteq\cdots\subseteq C(s_0)$. Transitivity of
set inclusion yields the claim.
\end{proof}

\begin{theorem}[Request-bounded effect theorem]
\label{thm:req-effect}
Let $q$ be a request-origin occurrence and $e$ an accepted protected effect in
the same ordinary authority segment. In a closed event-complete execution,
\begin{equation}
 \CausedBy(e,q)\land\Realize(e)=(o,r)
 \Longrightarrow (o,r)\in C(q).
 \label{eq:req-effect}
\end{equation}
\end{theorem}

\begin{proof}
$\CausedBy(e,q)$ implies $q\hb e$ by Theorem~\ref{thm:derived-ca}.
Lemma~\ref{lem:req-bound} then gives $C(e)\subseteq C(q)$. Since $e$ is
accepted and realizes $(o,r)$, Equation~\eqref{eq:effect-admission} gives
$(o,r)\in C(e)$. Therefore $(o,r)\in C(q)$.
\end{proof}

\subsection{Total model-level exclusion of the confused deputy}

\begin{definition}[Lineage-form confused deputy]
A confused-deputy state exists when there are a request-origin occurrence $q$,
an effect $e\in Eff$, an executor $D$, and a privilege $(o,r)$ such that
\begin{align}
  &\CausedBy(e,q),\\
  &\Realize(e)=(o,r),\\
  &\Exec(e)=D,\\
  &(o,r)\notin C(q),\\
  &\IndAuth(D,(o,r)).
  \label{eq:cda-def}
\end{align}
The final conjunct records the classical deputy situation: $D$ independently
possesses authority unavailable in the request context. It is not needed for
the contradiction below.
\end{definition}

\begin{theorem}[Total confused-deputy exclusion in the closed \Model model]
\label{thm:total-cda}
Every closed event-complete \Model execution satisfying ordinary PIC authority
continuity and protected-effect admission is free of lineage-form
confused-deputy states on every ordinary authority segment.
\end{theorem}

\begin{proof}
Assume a confused-deputy state satisfying \eqref{eq:cda-def}. By its first two
conjuncts and Theorem~\ref{thm:req-effect}, $(o,r)\in C(q)$. The definition
simultaneously requires $(o,r)\notin C(q)$, a contradiction. Independent
authority held by $D$ cannot enter the request's valid authority state through
an ordinary edge because \eqref{eq:authority-edge} forbids expansion.
\end{proof}

The result is ``total'' in the same formal sense in which RGR's Theorem~6 is
unconditional for FPs: every execution admitted by the stated semantics
satisfies the property. The two CDA definitions are not identified. RGR prove
their extensional $\CDAF$ property for their calculus; Theorem
\ref{thm:total-cda} proves the lineage-form origin-authority property for the
closed \Model event semantics. An explicit authorized endorsement or authority
relaxation is represented by a distinguished causal event and starts a new
segment; it is not a silent confused-deputy import of authority. This mirrors
RGR's treatment of endorsed regions as intentional, rather than confused, use
of lower-integrity influence.

The complete implication is
\begin{equation}
\boxed{
\begin{gathered}
\text{full-provenance influence on modeled events}\\
\Downarrow\\
\CausedBy(e,q)\Rightarrow q\hb e\\
\Downarrow\\
\text{PIC/COSS authority continuity}\\
\Downarrow\\
C(e)\subseteq C(q)\\
\Downarrow\\
\text{effect admission}\\
\Downarrow\\
\Realize(e)=(o,r)\Rightarrow(o,r)\in C(q)\\
\Downarrow\\
\neg\,\text{ConfusedDeputy}.
\end{gathered}}
\label{eq:full-cda-chain}
\end{equation}

\section{Re-reading the Capability/IFC Tension}

Consider the two-coordinate state
\begin{equation}
  \state=(C,F)
  \in X_{\mathrm{auth}}\times X_{\mathrm{ifc}}.
\end{equation}
A decision that reads only $C$ factors through the projection
\begin{equation}
  p_C(C,F)=C.
\end{equation}
If two states $(C,F_1)$ and $(C,F_2)$ require different information-flow
verdicts, any $C$-invariant mechanism must nevertheless treat them identically.
This is not a statement that capability systems are intrinsically incapable of
information-flow enforcement.  It is a statement about observational
insufficiency: the decision cannot enforce a distinction represented only in a
coordinate it does not observe.

This gives a compact reinterpretation of the historical sequence without
turning it into a contest between security paradigms.  Boebert's oracle does
observe clearance and classification when it grants a capability; the failure
in his construction occurs in subsequent propagation.  Even an oracle that
correctly creates capabilities cannot, in the unmodified machine he defines,
control their further propagation \cite{boebert1984}.  The distinction here is
therefore between the grant decision and what later executions may acquire and
exercise, not between a labeled and an unlabeled grant.

The later capability responses act directly on that propagation assumption.
Miller and Shapiro's data diode transmits data in one direction and capabilities
in neither \cite{miller_shapiro2003}; Miller, Yee, and Shapiro separately show
how distinguishing capability transfer from data transfer blocks the cross-level
capability movement required by Boebert's attack, while remaining within their
reading of his definition \cite{miller_yee_shapiro2003}.  Under a \Model
encoding, an oracle grant that enlarges the carried authority would be an
origination or explicitly authorized relaxation, whereas acquisition of a new
capability during ordinary continuation cannot silently enlarge the authority
coordinate.  This is a \Model interpretation of the security distinction, not
a claim that Boebert or the later capability papers used this state model.

Miller and Shapiro further show that program behavior and access abstractions can
tighten authority beyond what an arrangement-only analysis of the protection
graph predicts \cite{miller_shapiro2003}.  In \Model terms, these constructions
are examples of mechanisms that make additional distinctions observable or
restrict the semantic state on which the decision operates.  Once those
distinctions are observed, the mechanism no longer belongs to the
projection-blind class covered by the impossibility corollary.  The blindness
claim concerns the information read when an event is admitted; it does not say
that the original capability-granting oracle ignored clearance or
classification.

The same reasoning generalizes.  If
\begin{equation}
  \state=(S_1,S_2,\ldots,S_m),
\end{equation}
a policy that factors only through $S_1$ cannot enforce a property whose
verdict varies with $S_2$ while $S_1$ is fixed.  Security models can therefore
be individually valid and still be jointly necessary.

\section{PIC as a Derived Theory}

The relationship between \Model and PIC is derivational rather than
competitive.

\paragraph{PIC, first form.}
Take a single coordinate
\begin{equation}
  X=X_{\mathrm{auth}}=\powerset(O\times R).
\end{equation}
Then \Model-validity is exactly PoR plus authority attenuation, and Product
Continuity becomes PIC's origin-binding result.

\paragraph{Relational PIC.}
Take
\begin{equation}
  X=X_{\mathrm{auth}}\times X_{\mathrm{rel}}.
\end{equation}
Then the product order is
\begin{equation}
  (C',\rho')\preceq(C,\rho)
  \quad\Longleftrightarrow\quad
  C'\subseteq C\;\land\;\rho'\Rightarrow\rho,
\end{equation}
which is the product attenuation order already used by the relational theory
\cite{gallo_relational2026}.

\paragraph{General form.}
The natural closure is
\begin{equation}
  X=\prod_{j\in J}X_j,
\end{equation}
where authority, IFC, relational admissibility, provenance, geography,
separation-of-duty memory, or another security system may contribute a
coordinate.  PIC is therefore not renamed or invalidated.  It remains the
authority-continuity specialization from which the more general product
construction was discovered.

\section{COSS in the PIC Series}
\label{sec:series}

\Model is the third step in a sequence of models that share the same basic
shape: authorization is attached to an execution occurrence in causal history,
PoR relates adjacent occurrences, a carried state is ordered by refinement, and
ordinary continuation may preserve or restrict that state but not silently
expand it.  At the level of this continuity construction, each later paper
contains the earlier state discipline as a specialization; this statement does
not imply that the later construction reproduces every implementation theorem
or deployment assumption of the earlier papers.

\begin{table*}[t]
\centering
\small
\begin{tabular}{@{}p{0.15\linewidth}p{0.255\linewidth}p{0.255\linewidth}p{0.255\linewidth}@{}}
\toprule
 & PoC \cite{gallo_poc2026} & Relational PIC \cite{gallo_relational2026} & \Model \\
\midrule
Event domain & $O\times R\times L$ & execution occurrences indexed by causal histories & $\mathcal{B}\times L$ \\
Carried state & $C\subseteq O\times R$ & $(C,[\rho]_{\mathrm{sem}})$ & $\state\in\prod_j X_j$ \\
Order & $\subseteq$ & $\subseteq$ and $\Rightarrow$, componentwise & $\preceq_j$, componentwise \\
Continuity result & origin-bounded authority & Dual Monotonicity and Combined Closure & Product Continuity \\
Separation result & lineage-invariant policies cannot individuate occurrences & quotient factorization and observation-based indistinguishability & Projection Factorization and lineage-forgetting blindness \\
\bottomrule
\end{tabular}
\caption{The PIC series as a layered continuity construction.}
\label{tab:series}
\end{table*}

The containment statements are exact at the level claimed here.  With
$J=\{\mathrm{auth}\}$, the \Model transition rule is PoR plus
$C_{i+1}\subseteq C_i$, the validity rule of the first PoC model.  With
$J=\{\mathrm{auth},\mathrm{rel}\}$ and relational refinement ordered by
implication, the product order is the relational PIC attenuation order
$(C',\rho')\preceq(C,\rho)$ iff $C'\subseteq C$ and
$\rho'\Rightarrow\rho$ \cite{gallo_relational2026}.

The separation results share the same factorization principle but should not be
identified more strongly than their domains permit.  For the authority-only
event space, take $b=(o,r)$ and two histories $\ell,\ell'$ such that
$(o,r)\in C(\ell)$ and $(o,r)\notin C(\ell')$.  The \Model
lineage-forgetting corollary then gives the same required separation as the
first paper's lineage-invariance result.  The relational paper subsequently
shows that the earlier lineage-forgetting projection is one quotient instance
of a more general factorization construction; its executor-profile quotients and
\Model's coordinate projections are applications of the same mathematical
principle, not literally the same projection map \cite{gallo_relational2026}.

The confused-deputy result is strengthened in the present paper. Hardy's
classical example motivates the problem of a deputy exercising authority for
the wrong request context \cite{hardy1988}. The first PoC paper excludes
authority import as valid behavior inside a valid lineage through origin-bounded
attenuation. The relational paper then isolates Complete Mediation and Causal
Attribution as the premises needed to lift that closure to actual request-caused
effects \cite{gallo_relational2026}. Section~\ref{sec:fp-axis} internalizes
those premises in a closed event semantics: protected effects are modeled
admission events, and RGR-style full-provenance reachability derives causal
attribution on the same occurrence order as PoR. The resulting
Theorem~\ref{thm:total-cda} excludes every lineage-form confused-deputy state
along ordinary continuation in the closed model, while preserving explicit endorsement/authorized
relaxation as a distinct causal transition class.

\section{Heterogeneous Security-State Spaces}

PIC already permits policy translation between heterogeneous operation/resource
spaces \cite{gallo_poc2026}.  The same idea extends componentwise.

At hop $i$, let coordinate $j$ use state space $X_i^{(j)}$ with order
$\preceq_i^{(j)}$.  Let
\begin{equation}
  T_{i\rightarrow i+1}^{(j)}:
  X_i^{(j)}\to X_{i+1}^{(j)}
\end{equation}
be an order-theoretic policy translation.  A translated transition satisfies
\begin{equation}
  \coord_{i+1}^{(j)}
  \preceq_{i+1}^{(j)}
  T_{i\rightarrow i+1}^{(j)}(\coord_i^{(j)}).
  \label{eq:translated-transition}
\end{equation}

If the translations are monotone and composable, define
\begin{equation}
  T_{0\rightarrow n}^{(j)}
  =T_{n-1\rightarrow n}^{(j)}\circ\cdots\circ T_{0\rightarrow1}^{(j)}.
\end{equation}

\begin{proposition}[Translated continuity]\label{prop:translated}
If every $T^{(j)}_{i\rightarrow i+1}$ is monotone and every transition
satisfies \eqref{eq:translated-transition}, then
\begin{equation}
  \coord_n^{(j)}\preceq_n^{(j)}T_{0\rightarrow n}^{(j)}(\coord_0^{(j)}).
\end{equation}
\end{proposition}

\begin{proof}
By induction on $n$. For $n=0$ the claim is reflexivity. Suppose
$\coord_i^{(j)}\preceq_i^{(j)}T_{0\rightarrow i}^{(j)}(\coord_0^{(j)})$.
Monotonicity of $T^{(j)}_{i\rightarrow i+1}$ gives
$T^{(j)}_{i\rightarrow i+1}(\coord_i^{(j)})\preceq_{i+1}^{(j)}
T_{0\rightarrow i+1}^{(j)}(\coord_0^{(j)})$, and
\eqref{eq:translated-transition} with transitivity gives the claim for
$i+1$.
\end{proof}

Proposition~\ref{prop:translated} is order-theoretic: it requires monotonicity
of the translations and proves a bound in the target orders. A stronger claim
that a translation preserves the intended security meaning across heterogeneous
behavior universes requires a separate semantic relation between those
universes. This paper does not define such a cross-domain semantic soundness
condition; an overly permissive semantic translation is therefore a policy or
deployment error rather than a failure of the order-theoretic continuity
theorem.

\section{Ordinary Continuation versus Authorized Relaxation}
\label{sec:relaxation}

Not every legitimate security transition is monotone.  Declassification may
relax an IFC constraint; a new authority grant may enlarge an authority state;
a policy administrator may broaden an admissible geographic region.

\Model therefore distinguishes two classes of transition.

\begin{definition}[Ordinary continuation]
An ordinary continuation is a PoR-linked transition satisfying
$\state_{i+1}\preceq\state_i$ (or the translated form of
\eqref{eq:translated-transition}).
\end{definition}

\begin{definition}[Authorized relaxation event]
An authorized relaxation event is an explicitly distinguished causal event
whose validity is justified by a separate authority or policy rule and which
may replace one or more coordinates by states not below their predecessors.
\end{definition}

The point is not to prohibit declassification, re-authorization, or policy
change.  It is to prevent those events from being confused with ordinary
continuation.  A security state may become more permissive only at a boundary
whose authority to perform that relaxation is itself explicit and auditable.

This mirrors the distinction already present in PIC between attenuation within
one lineage and the creation of new authority at an origin or explicit future
composition boundary.

Relaxation events are still causal events: they are PoR-linked to their
predecessor, and only the monotonicity conjunct is replaced by the separate
authorization rule.

\begin{proposition}[Segment continuity]\label{prop:segment}
Let a PoR-linked lineage contain authorized relaxation events at indices
$1\le r_1<\cdots<r_k$, where the event at $r_m$ is the transition from
$s_{r_m-1}$ to $s_{r_m}$, all other transitions being ordinary
continuations. For every $i$, let $r(i)$ be the largest $r_m\le i$, or $0$
if none exists. Then
\begin{equation}
  \state_i\preceq\state_{r(i)}.
\end{equation}
\end{proposition}

The statement assumes a homogeneous state space along the segment. For segments
containing translated hops, the corresponding bound is the one given by
Proposition~\ref{prop:translated}.

\begin{proof}
If $r(i)=i$, the claim is reflexivity. Otherwise, every transition leaving
$s_{r(i)}$ and ending at or before $s_i$ is an ordinary continuation, because
$r(i)$ is the most recent relaxation-entry index. Product Continuity applied
to the segment from $s_{r(i)}$ through $s_i$ gives
$\state_i\preceq\state_{r(i)}$.
\end{proof}

Origin-boundedness is therefore replaced, in the presence of relaxation,
by boundedness from the most recent explicitly authorized relaxation. A
relaxation cannot be introduced by ordinary continuation; it can only
restart the bound at an explicit, auditable point.

\section{Causal DAGs, Joins, and Multiple Causes}

The common causal substrate is a partial order, not inherently a linear call
chain. Linear lineage notation denotes a path in that order. Fan-out, fan-in,
shared state, asynchronous delivery, and multiple-cause effects are therefore
represented by DAG structure before any security coordinate is considered.

For an ordinary event $s$ with immediate predecessors
$\mathrm{Pred}(s)=\{p_1,\ldots,p_k\}$, the conservative \Model rule is
\begin{equation}
  \st(s)\preceq \st(p_m)
  \qquad\text{for every }m\in\{1,\ldots,k\}.
  \label{eq:dag-join}
\end{equation}
For the authority coordinate this specializes to
\begin{equation}
  C(s)\subseteq\bigcap_{m=1}^{k}C(p_m).
\end{equation}
Equation~\eqref{eq:dag-join} is sufficient for Product Continuity along every
parent path and for the request-bounded effect theorem even when an effect has
several causes. In particular, if $q$ is any cause whose influence reaches
$s$, then the authority exercised at $s$ remains bounded by $C(q)$ unless an
explicit authorized relaxation/composition event intervenes.

A non-conservative composition policy may intentionally combine authority or
other security states from several parents. Such a transition is not an
ordinary continuation and is represented as an explicit authorized composition
or relaxation event. Its policy semantics must state which origins justify the
new state. This keeps the general theory honest without treating queues,
databases, timers, callbacks, or fan-in as failures of causal attribution: those
mechanisms contribute causal edges; only a genuine broadening of carried
security state requires a separate authorization rule.

Miller and Shapiro's ``mutually suspicious composition'' concerns the
coexistence of independently motivated access abstractions over one capability
graph \cite{miller_shapiro2003}. The connection is conceptual rather than an
identity of formal operators: \Model's DAG rule is an event-level condition on
how the security states of multiple causal parents may continue or be
explicitly recomposed.

\section{Related Work}

\paragraph{Bell--LaPadula.}
Bell and LaPadula model secure system states together with rules of operation
that preserve security under state transitions \cite{bell_lapadula1973model}.
\Model does not replace those rules or claim that Bell--LaPadula is merely a
set-inclusion policy.  Rather, an effective Bell--LaPadula security state may
serve as one coordinate when its local refinement relation and admitted-event
semantics are specified; \Model adds the explicit causal-event axis and the
product transport rule used when that state must coexist with other security
states.

\paragraph{Denning's lattice model.}
Denning organizes security classes as a lattice whose order is justified by the
semantics of permitted information flow \cite{denning1976}.  This is a direct
example of an ordered security domain, but the IFC semantics remain Denning's:
\Model supplies neither the lattice nor a program analysis for discovering
flows.  It supplies a way to carry such an ordered state together with other
coordinates along one causal execution.

\paragraph{Execution-history access control.}
Abadi and Fournet make run-time rights depend on execution history rather than
only the current stack \cite{abadi_fournet2003}.  In their model, ordinary
execution automatically intersects current rights with the static rights of
code that runs, while an explicit special operation may modify current rights,
at most restoring the static rights of that code and subject to further
constraints.  This is structurally close to the distinction here
between ordinary non-expansive continuation and an explicit relaxation event.
The models nevertheless have different domains: their construction computes a
rights state for code execution history, whereas \Model abstracts over an
arbitrary finite product of security-state coordinates and takes sound PoR as
an explicit premise relating the modeled execution occurrences.

\paragraph{Decentralized information-flow control.}
Myers and Liskov define decentralized labels with owner-specific reader sets,
a restriction order that forms a security-class lattice, and explicit
declassification that is legal only when the process has authority to act for
the relevant owner \cite{myers_liskov1997}.  Their restriction/declassification
distinction is therefore an important prior instance of monotone security
restriction plus explicitly authorized relaxation.  \Model treats such an IFC
state as one possible coordinate; it does not replace the label semantics,
acts-for relation, or static flow analysis of that work.

\paragraph{Rajani--Garg--Rezk full provenance.}
RGR give the first language-based extensional characterization of CDA-freedom,
show that capability semantics does not prevent all implicit-designation CDAs,
and introduce explicit-only and full provenance semantics
\cite{rajani_garg_rezk2016}. FPs augments value provenance with a
program-counter label stack so that branch and loop conditions contribute to
the integrity check on later writes. Their Theorem~6 proves CDA-freedom with
endorsement unconditionally for their region calculus, via information-flow
integrity and their Theorem~7. Read from \Model, FPs is a language-level
instance of the influence relation required by Section~\ref{sec:fp-axis}.
\Model does not reinterpret their theorem as a protocol result; the new step is
a conservative event refinement that places FPs-style influence on the same
execution occurrences as PoR and uses that common order to derive causal
attribution for the lineage-form authority theorem.

\paragraph{Unified authorization and dependency calculi.}
Flow-limited authorization integrates authorization and information-flow
reasoning for delegation and trust \cite{arden_liu_myers2015}; the Dependency
Core Calculus gives one calculus in which several dependency analyses,
including information flow, are instances \cite{abadi_dcc1999}; and
operating-system DIFC designs track labels and privileges across processes
\cite{asbestos2005,histar2006,flume2007}. These works unify security reasoning
within a language semantics or an operating-system label model. \Model starts
from a different object: authority bound to execution occurrences in a causal
history that may span services and trust domains, with continuity verifiable
at the receiving boundary. The approaches are complementary: a language-based
or operating-system analysis can supply the influence relation, or a security
coordinate, that \Model transports along PoR.

\paragraph{Capability systems and access abstractions.}
Miller, Yee, and Shapiro show that conclusions about capability systems depend
on the capability model and on distinctions such as data transfer versus
capability transfer, designation with authority, and access-controlled
delegation channels \cite{miller_yee_shapiro2003}.  Miller and Shapiro further
show that security-enforcing program behavior and access abstractions can
restrict authority beyond what a protection-graph-only analysis exposes
\cite{miller_shapiro2003}.  \Model's claim is not that these systems lack
context, flow control, or causal reasoning.  Its contribution is a uniform
event/state construction and a factorization result: whenever a target verdict
depends on a distinction erased by an observation, a decision invariant under
that observation cannot reproduce the target verdict.

\section{Semantic Closure and Realization Assumptions}

The closed formal model now contains the two facts that earlier PIC work left as
separate deployment assumptions. First, every protected effect is itself a
modeled occurrence to which product admission is applied. Second, every
security-relevant explicit data/designation dependency and implicit control
dependency in the stated threat model is represented by the full-provenance
influence relation. Theorem~\ref{thm:derived-ca} then derives causal attribution
from reachability.

This change is important for asynchronous systems. A queue, database, shared
file, cache entry, timer, callback, or message broker does not break the causal
chain merely because the original stack frame has ended. If the later event
consumes state produced by the earlier event, the write/read, enqueue/dequeue,
schedule/fire, or send/receive dependency is an edge of the same event DAG.
The security state transported by \Model is therefore attached to causal
occurrences rather than to synchronous call-stack lifetime.

A concrete realization must still justify that its evidence faithfully
implements the formal semantics. In particular:
\begin{enumerate}[leftmargin=*]
  \item PoR evidence must be sound for the modeled immediate causal edges;
  \item the runtime or analysis providing full provenance must conservatively
        expose every influence class claimed by the threat model;
  \item each security coordinate and heterogeneous translation must faithfully
        represent the security meaning assigned to its refinement order; and
  \item explicit endorsement, declassification, authority grant, or composition
        events must be separately authorized and observable as such.
\end{enumerate}
These are realization obligations, not an additional Causal Attribution axiom.
If a physical influence channel is omitted from the formal event semantics, no
theorem about the model can silently cover that omitted channel.

\section{Limitations and Open Questions}

The causal-attribution item present in the earlier draft is discharged for the
closed event model of Section~\ref{sec:fp-axis}. The remaining research tasks
are narrower.

\paragraph{Semantic soundness of heterogeneous translations.}
Translated Continuity proves preservation relative to the supplied ordered
translation maps. Relating states and behaviors across different local
vocabularies still requires a cross-domain semantic relation and a soundness
condition for each translation.

\paragraph{Order soundness and coordinate choice.}
The meaning of $\preceq_j$ is domain-specific. A mathematically monotone but
semantically wrong order proves the wrong property. Likewise, omitting a
coordinate is sound only for target properties invariant under that omission;
the projection theorem makes this requirement explicit.

\paragraph{Physical coverage.}
The formal theorem is total over the events and influence channels represented
by the model. A real system must establish that the threat model's relevant
physical channels are represented by the execution/provenance semantics.
Covert or side channels not represented as causal influence remain outside the
formal domain rather than counterexamples to the theorem.

\paragraph{Implementation evidence.}
The theory remains agnostic about whether evidence is carried in tokens,
attestations, trusted-runtime state, gateways, proof-carrying objects, dynamic
information-flow tracking, or another mechanism. Representation is secondary
to whether the mechanism soundly establishes the common causal order and the
security coordinates needed by the property being claimed.

\section{Discussion}

\paragraph{Scope of the framework.}
Across the three-paper sequence, the common primitives are an execution
occurrence, a causal history, PoR between adjacent occurrences, an ordered
state carried by the history, ordinary continuation that does not expand that
state, and an explicit transition class for authorized relaxation.  The
continuity and separation results arise from these primitives plus the stated
semantic and enforcement assumptions.  The order theory used here is
elementary; the contribution claimed is the choice of authorization domain and
the causal transport discipline, not a new result in order theory.

The proposed shift is small syntactically but large semantically. Instead of
asking one security model to become universal, \Model treats each model as a
coordinate or influence view over the same causal execution. The RGR
integration makes the distinction concrete: the IFC policy state is one
coordinate, while full provenance is an influence suborder on the same
occurrences. Neither requires a second temporal axis.

The general invariant is:
\begin{quote}
Ordinary causal transport may not make any carried security state more
permissive according to that security system's own refinement order.  Any
relaxation must be an explicit, separately authorized causal event.
\end{quote}

This separates three concerns that are often conflated:
\begin{enumerate}[leftmargin=*]
  \item \emph{what} a security system considers more or less permissive;
  \item \emph{which} security systems a decision must observe; and
  \item \emph{whether} the observed state belongs to the causal execution that
        is actually producing the effect.
\end{enumerate}

Capability authority, Bell--LaPadula-style information-flow constraints,
relational executor admissibility, provenance, geography, and other controls
can therefore remain distinctly meaningful.  They need not be merged into a
single ontology.  Their product is carried by the execution, while PoR provides
the causal axis that keeps the evolving states attached to one history.

An access abstraction need not always become a new coordinate.  It may instead
be the concrete mechanism that realizes a coordinate's semantic extension, or
it may restrict the admissible subset $Y$ of product states.  This is important
for reading the capability literature correctly: the factory, Caretaker, and
data-diode constructions of Miller and Shapiro constrain authority by program
behavior, even when a coarser projection of the protection state would admit
more possibilities \cite{miller_shapiro2003}.  \Model abstracts over the
implementation choice and asks only which distinctions the receiving decision
actually observes and preserves along the causal execution.

\section{Conclusion}

PIC placed authority-use events on a causal axis and made authority continuation
non-expansive. \Model keeps that axis and generalizes the carried state from one
authority set to a product of heterogeneous security states. The Product
Continuity theorem is componentwise transitivity: ordinary causal transport
cannot introduce a behavior forbidden by an origin state in any carried
coordinate. The projection theorem gives the dual limitation: a decision
cannot enforce a distinction that its observation erases.

Read from PIC, RGR provide an independent, language-level treatment of the
confused deputy. Their
capability semantics blocks explicit designation but not every implicit attack;
their EP semantics tracks explicit provenance; and their FPs semantics adds
program-counter provenance for implicit control flow. Their Theorem~6 proves
CDA-freedom with endorsement unconditionally for their complete region calculus
\cite{rajani_garg_rezk2016}. The present paper does not rewrite that theorem as
PoR or PoC. Instead it lifts the influence represented by FPs to ghost event
dependencies whose erasure leaves the original FPs execution unchanged.

This event lift resolves the causal question at the right level. A modeled
machine execution has one causal order of occurrences. Authority, IFC policy,
relational admissibility, and other security states are attached to those
occurrences, while full provenance identifies which earlier occurrences
influence later data or control. Full-provenance influence is therefore a
suborder of the same execution order. Asynchronous state transfer through a
database, queue, timer, callback, shared object, or message does not create a
new time; it contributes additional causal dependency edges.

Once the modeled execution is closed under protected effects and full
provenance, causal attribution is reachability. Combining that fact with
ordinary PIC authority attenuation gives
\begin{equation}
  \CausedBy(e,q)\land\Realize(e)=(o,r)
  \Longrightarrow (o,r)\in C(q).
\end{equation}
The lineage-form confused-deputy condition requires the opposite,
$(o,r)\notin C(q)$, and is therefore unsatisfiable along ordinary
continuation. Independent authority held
by the deputy cannot be imported through an ordinary causal continuation;
intentional endorsement, authority grant, or non-conservative composition is a
distinguished authorized event and begins a new authority bound.

Accordingly, the final model-level result is
\begin{equation}
\boxed{
\begin{gathered}
\text{closed event-complete full provenance}\\
{}+\ \text{COSS/PIC authority continuity}\\
{}+\ \text{protected-effect admission}\\
\Longrightarrow\ \text{confused-deputy exclusion along ordinary continuation}.
\end{gathered}}
\end{equation}
``Total'' has the same formal reading as any semantic safety theorem: no
execution admitted by the model contains the forbidden state. It does not
claim security for physical influence channels omitted from the model. Within
the stated event semantics, however, Causal Attribution is no longer an
independent assumption; it is a derived theorem of the common causal order.

\begin{thebibliography}{99}

\bibitem{bell_lapadula1973}
D.~E. Bell and L.~J. LaPadula,
``Secure Computer Systems: Mathematical Foundations,''
The MITRE Corporation, ESD-TR-73-278, Vol.~I, 1973.

\bibitem{bell_lapadula1973model}
L.~J. LaPadula and D.~E. Bell,
``Secure Computer Systems: A Mathematical Model,''
The MITRE Corporation, ESD-TR-73-278, Vol.~II, 1973.

\bibitem{denning1976}
D.~E. Denning,
``A Lattice Model of Secure Information Flow,''
\emph{Communications of the ACM}, vol.~19, no.~5, pp.~236--243, 1976,
doi: 10.1145/360051.360056.

\bibitem{boebert1984}
W.~E. Boebert,
``On the Inability of an Unmodified Capability Machine to Enforce the *-Property,''
in \emph{Proceedings of the 7th DoD/NBS Computer Security Conference},
Gaithersburg, MD, USA, Sep.~24--26, 1984, pp.~291--293.
[Online]. Available: \url{https://csrc.nist.gov/files/pubs/conference/1984/09/24/7th-dodnbs-computer-security-conference/final/docs/1984-7th-conference-proceedings.pdf}.

\bibitem{miller_yee_shapiro2003}
M.~S. Miller, K.-P. Yee, and J.~S. Shapiro,
``Capability Myths Demolished,''
Technical Report SRL2003-02, Systems Research Laboratory,
Department of Computer Science, Johns Hopkins University, Mar. 2003.

\bibitem{miller_shapiro2003}
M.~S. Miller and J.~S. Shapiro,
``Paradigm Regained: Abstraction Mechanisms for Access Control,''
in \emph{Advances in Computing Science -- ASIAN 2003: Programming Languages and Distributed Computation},
Lecture Notes in Computer Science, vol.~2896, Springer, 2003,
pp.~224--242, doi: 10.1007/978-3-540-40965-6\_15.

\bibitem{abadi_fournet2003}
M.~Abadi and C.~Fournet,
``Access Control Based on Execution History,''
in \emph{Proceedings of the 10th Annual Network and Distributed System Security Symposium (NDSS)},
2003.

\bibitem{myers_liskov1997}
A.~C. Myers and B.~Liskov,
``A Decentralized Model for Information Flow Control,''
in \emph{Proceedings of the 16th ACM Symposium on Operating Systems Principles (SOSP)},
pp.~129--142, 1997, doi: 10.1145/268998.266669.

\bibitem{rajani_garg_rezk2016}
V.~Rajani, D.~Garg, and T.~Rezk,
``On Access Control, Capabilities, Their Equivalence, and Confused Deputy Attacks,''
in \emph{2016 IEEE 29th Computer Security Foundations Symposium (CSF)},
Lisbon, Portugal, 2016, doi: 10.1109/CSF.2016.18.

\bibitem{arden_liu_myers2015}
O.~Arden, J.~Liu, and A.~C. Myers,
``Flow-Limited Authorization,''
in \emph{2015 IEEE 28th Computer Security Foundations Symposium (CSF)}, 2015.

\bibitem{abadi_dcc1999}
M.~Abadi, A.~Banerjee, N.~Heintze, and J.~G. Riecke,
``A Core Calculus of Dependency,''
in \emph{Proceedings of the 26th ACM SIGPLAN-SIGACT Symposium on Principles of
Programming Languages (POPL)}, 1999.

\bibitem{asbestos2005}
P.~Efstathopoulos, M.~Krohn, S.~VanDeBogart, C.~Frey, D.~Ziegler, E.~Kohler,
D.~Mazi\`eres, F.~Kaashoek, and R.~Morris,
``Labels and Event Processes in the Asbestos Operating System,''
in \emph{Proceedings of the 20th ACM Symposium on Operating Systems Principles
(SOSP)}, 2005.

\bibitem{histar2006}
N.~Zeldovich, S.~Boyd-Wickizer, E.~Kohler, and D.~Mazi\`eres,
``Making Information Flow Explicit in HiStar,''
in \emph{Proceedings of the 7th USENIX Symposium on Operating Systems Design and
Implementation (OSDI)}, 2006.

\bibitem{flume2007}
M.~Krohn, A.~Yip, M.~Brodsky, N.~Cliffer, M.~F. Kaashoek, E.~Kohler, and
R.~Morris,
``Information Flow Control for Standard OS Abstractions,''
in \emph{Proceedings of the 21st ACM Symposium on Operating Systems Principles
(SOSP)}, 2007.

\bibitem{lamport1978}
L.~Lamport,
``Time, Clocks, and the Ordering of Events in a Distributed System,''
\emph{Communications of the ACM}, vol.~21, no.~7, pp.~558--565, 1978,
doi: 10.1145/359545.359563.

\bibitem{hardy1988}
N.~Hardy,
``The Confused Deputy (or Why Capabilities Might Have Been Invented),''
\emph{ACM SIGOPS Operating Systems Review}, vol.~22, no.~4, pp.~36--38, 1988.

\bibitem{gallo_poc2026}
N.~Gallo,
``Proof-of-Continuity: A Temporal Model for Authority Propagation in Distributed Systems and AI Agents,''
arXiv:2607.08906, 2026.

\bibitem{gallo_relational2026}
N.~Gallo,
``Beyond Proof of Possession: A Relational Theory of Authority,''
manuscript, 2026.

\end{thebibliography}

\end{document}
